\begin{claim}{D.2.53}{Proved}{The epsilon-prime target term is zero}
Let $X_0=\Sigma K(\FF_2,2)$, $Y=\Sigma X_0$, and
$u_3:S^3\to X_0$, $u_4=Eu_3:S^4\to Y$ be the bottom maps.
Then
\[
                         u_4E\epsilon'=0\quad\text{in }\pi_{14}Y.
                         \tag{D.2.53a}
\]
More precisely, the factor $t$ in D.2.49 can be chosen with
the compatible Toda extension and nullhomotopies so that
\[
 u_3\epsilon'=\pm t\nu_{10},\qquad
 t\in\{u_3,\nu',2\nu_6\},\qquad Et=0.
 \tag{D.2.53b}
\]
This identifies a chosen factorization; it does not assert
that every map $t$ satisfying the first equality is zero
after suspension.

Consequently the generator and the actual Moore obstruction
of D.2.46 have the sharper form
\[
 \begin{gathered}
 q=\alpha\sigma',\qquad 4q=0,\qquad
 h=h_Ml_4=n\alpha\sigma'\text{ for some integer }n,\\
                         2h=2\alpha\sigma'.
 \end{gathered}
 \tag{D.2.53c}
\]
The same zero holds for the unique lift of the term in
\textup{(D.2.53a)} to $L=Y\langle7\rangle$.
Neither the value nor the order of $\alpha\sigma'$ in
the full target is determined.  If it has order four,
$h=\pm\alpha\sigma'$; if it has order at most two,
$h$ may be zero or $\alpha\sigma'$.
\end{claim}
\begin{proof}
All Toda brackets below are ordinary brackets in the
two-local pointed homotopy category.  The actual homotopy
groups of $Y$ in question are two-primary, so this does
not lose any part of the target assertion.
We use the actual identities
\[
             u_3\nu'=0,\qquad u_4E\nu'=0,
                         \qquad 2\pi_{11}Y=0,
             \tag{D.2.53d}
\]
proved in D.2.49, D.2.44 and D.2.31, respectively.

\textbf{A compatible representative, not an arbitrary factor.}
The indexed definition of $\epsilon'$ used in D.2.49 lies
in the ordinary bracket
\[
                    \epsilon'\in\{\nu',2\nu_6,\nu_9\}.
                    \tag{D.2.53e}
\]
This is the definition and indexed-bracket convention in
\cite[Section 2.5, p.~4, and the accompanying author program]
{FGHStricklandToda}; the original Toda attribution is
Section VI(iv).  Choose an extension and a coextension
realizing this specified value:
\[
 \begin{gathered}
 C=C(2\nu_6)=S^6\cup_{2\nu_6}e^{10},\qquad
 \bar\nu':C\longrightarrow S^3,\qquad
 \widetilde\nu_9:S^{13}\longrightarrow C,\\
 \bar\nu'|_{S^6}=\nu',\qquad
 \bar\nu'\widetilde\nu_9=\epsilon',\qquad
 p\widetilde\nu_9=\nu_{10},
 \quad p:C\longrightarrow S^{10}.
 \end{gathered}
 \tag{D.2.53f}
\]
The harmless common sign depends on the cone convention.
Use a nullhomotopy of $u_3\nu'$ in
\textup{(D.2.53d)} to factor the actual composite as
\[
                         u_3\bar\nu'\simeq t p.
                         \tag{D.2.53g}
\]
The factor obtained from these two compatible
nullhomotopies is precisely, up to that common sign,
a value of $\{u_3,\nu',2\nu_6\}$.  This is the
extension--coextension definition of the ordinary bracket:
the nullhomotopy of $\nu'(2\nu_6)$ defining $\bar\nu'$
and the chosen nullhomotopy of $u_3\nu'$ are glued along
their common face.  Thus this conclusion does not follow
merely from the existence of some pinch factor in a
Puppe sequence.  Composing \textup{(D.2.53g)} with the
\emph{same} $\widetilde\nu_9$ gives
$u_3\epsilon'=\pm t\nu_{10}$.
Suspension of this chosen cone diagram gives
\[
                    Et\in\mathcal B_2
                         :=\{u_4,E\nu',2\nu_7\}
                              \subseteq\pi_{11}Y.
                    \tag{D.2.53h}
\]
No inverse suspension or identification of unrelated
nullhomotopy choices is used.

\textbf{The suspended bracket is the singleton zero.}
The exact two-primary sphere groups needed here are
\[
                 \pi_9(S^3)_{(2)}=0,\qquad
                 \pi_{10}(S^3)_{(2)}=0,\qquad
                 \pi_{11}(S^4)_{(2)}=0.
                 \tag{D.2.53i}
\]
They are explicitly tabulated in the opened author
program accompanying \cite{FGHStricklandToda}, in its
sixth- and seventh-stem blocks, with the attributions
to Toda Proposition 5.11, p.~46, and Proposition 5.15,
p.~50.  The last equality also follows from the first
relevant terms of the Hopf-fibration sequence
$S^3\to S^7\to S^4$, since
$\pi_{11}S^7=\pi_4^s\mathbb S=0$ and
$\pi_{10}(S^3)_{(2)}=0$.

Since $\nu_7=E\nu_6$, the first equality in
\textup{(D.2.53i)} gives
\[
                    (E\nu')\nu_7=E(\nu'\nu_6)=0.
\]
Hence the smaller bracket
$\mathcal B_1=\{u_4,E\nu',\nu_7\}$ is defined.
Precompose its last sphere by the degree-two map
$2\iota_{10}$.  Naturality of the ordinary cone diagram
gives
\[
                    \mathcal B_1\circ(2\iota_{11})
                               \subseteq\mathcal B_2.
                    \tag{D.2.53j}
\]
This is precomposition on a sphere, not an assumed
additivity formula for the middle bracket entry.
Every value on the left is zero by $2\pi_{11}Y=0$.
Thus $0\in\mathcal B_2$.

The complete ordinary indeterminacy of $\mathcal B_2$ is
\[
                 u_4\circ\pi_{11}S^4
                     +\pi_8Y\circ(2\nu_8).
                 \tag{D.2.53k}
\]
The first term is zero by \textup{(D.2.53i)}.
For every $\gamma:S^8\to Y$, the second term is zero
by the exact right-composition identity
\[
              \gamma(2\nu_8)
                  =\gamma\nu_8(2\iota_{11})
                  =2(\gamma\nu_8)=0.
\]
Again only the exponent of $\pi_{11}Y$ is used; no
distributivity in the variable $\gamma$ is needed.
The ordinary bracket is a coset of this zero
indeterminacy, so
\[
                            \mathcal B_2=\{0\}.
                            \tag{D.2.53l}
\]
Equations \textup{(D.2.53h)--(D.2.53l)} prove $Et=0$
for the compatible factor selected above.  Suspending
$u_3\epsilon'=\pm t\nu_{10}$ now proves
$u_4E\epsilon'=\pm(Et)\nu_{11}=0$.

The connected-cover map $L\to Y$ is an isomorphism on
homotopy groups above degree six; the unique lifted
term is consequently zero as well.  Substituting this
value into D.2.46 proves \textup{(D.2.53c)} and its
two order-dependent alternatives.  It does not
evaluate the remaining $\alpha\sigma'$ composite.
\end{proof}
